\section{Blockchain in Finance: The Competitive Landscape}

\subsection{The Cross-Border Payment Problem \& Bitcoin's Solution}
The legacy cross-border payment system (built in 1973) relies on correspondent
banking, suffering from high costs ($\sim 6.25\%$), slow settlement (1-5
business days), and the capital inefficiency of pre-funded Nostro/Vostro
accounts.
\begin{itemize}
    \item \textbf{Bitcoin's Solution:} Bitcoin solves the problem of reaching consensus among
          untrusted strangers (the Byzantine Generals Problem) by replacing a central
          bookkeeper with a decentralized, cryptographically secured ledger.
    \item \textbf{Key Mechanisms:}
          \begin{itemize}
              \item \textit{Identity \& Authorization:} Public/private key cryptography and digital
                    signatures.
              \item \textit{Immutability \& Ordering:} Each block contains the hash of the previous block.
              \item \textit{Consensus (Mining):} Proof of Work (PoW) prevents Sybil attacks by requiring
                    computational effort.
              \item \textit{Double Spend Prevention:} Wait for 6 block confirmations.
          \end{itemize}
    \item \textbf{The Trade-off:} Bitcoin achieves permissionless, censorship-resistant,
          trustless consensus but fails as an institutional payment rail due to low
          throughput (3.3-7 TPS), massive energy consumption (1000+ kWh/tx), price
          volatility, and pseudonymity (incompatible with KYC/AML).
\end{itemize}

\subsection{Permissionless vs. Permissioned Blockchains}
\begin{itemize}
    \item \textbf{Permissionless (e.g., Bitcoin, Ethereum):} Anyone can join/validate. Validators
          are anonymous (miners/stakers). High decentralization but slow and lacks
          KYC/AML.
    \item \textbf{Permissioned (e.g., Ripple/XRPL, JPMorgan Kinexys):} Approved participants
          only. Known, vetted validators. High speed, clear KYC/AML, but sacrifices pure
          decentralization.
\end{itemize}

\subsection{The BIS Framework for System-Level Money}
According to the Bank for International Settlements (BIS), any instrument
claiming to function as system-level money must pass three tests:
\begin{enumerate}
    \item \textbf{Singleness:} Can it be accepted at par across the system? (1 unit always = 1
          unit, regardless of issuer).
    \item \textbf{Elasticity:} Can supply expand to meet payment demand (like credit lines)
          without requiring 100\% pre-funding?
    \item \textbf{Integrity:} Is there a functional, system-level KYC/AML layer to resist
          financial crime?
\end{enumerate}

\subsection{Stablecoins, CBDCs, and XRP}
\begin{itemize}
    \item \textbf{Stablecoins:} Private-sector digital tokens pegged to fiat (e.g., USDT, USDC).
          \begin{itemize}
              \item \textit{BIS Verdict:} Fails Singleness (tokens are issuer-tagged), fails Elasticity
                    (requires 100\% pre-funding, imposing a cash-in-advance constraint), and fails
                    Integrity (pseudonymous bearer instruments).
          \end{itemize}
    \item \textbf{CBDCs (Central Bank Digital Currencies):} Digital liabilities of the central
          bank.
          \begin{itemize}
              \item \textit{Retail CBDC:} Issued to the public (e.g., China's e-CNY). Risks
                    disintermediating commercial banks.
              \item \textit{Wholesale CBDC:} Issued only to financial institutions for interbank settlement
                    (e.g., Project mBridge). Directly competes with correspondent banking.
          \end{itemize}
    \item \textbf{XRP (Ripple):} A bridge currency for cross-border settlement. It is neither a
          stablecoin (price fluctuates), nor a CBDC, nor a decentralized permissionless
          crypto. It operates on a private/permissioned network to eliminate
          Nostro/Vostro pre-funding via instant (3-5 second) settlement.
\end{itemize}

\subsection{The Regulatory Frame: Who Gets to Play}
Regulatory jurisdiction acts as a strategic variable that shapes the
competitive landscape:
\begin{itemize}
    \item \textbf{US GENIUS Act (2025):} Strictly regulates stablecoins. Requires a 1:1 reserve
          in liquid US assets (dollars, T-bills < 93 days), mandates monthly public
          disclosures, prohibits paying yield/interest, and creates severe hurdles for
          non-US issuers (like Tether), giving a structural advantage to US-domiciled
          issuers like USDC.
    \item \textbf{Singapore PSA (MAS):} An activity-based, principles-driven framework. Requires
          a Major Payment Institution (MPI) license for cross-border transfers and
          mandates 1:1 reserves for Single-Currency Stablecoins with strict Travel Rule
          enforcement.
    \item \textbf{The SEC vs. Ripple Case (The Howey Test):} The SEC sued Ripple alleging XRP was
          an unregistered security.
          \begin{itemize}
              \item \textit{2023 Ruling:} XRP sold on public retail exchanges is \textbf{not} a security. XRP sold directly to institutional investors \textbf{is} a security.
              \item \textit{2025 Resolution:} Settled with civil penalties; a permanent injunction on direct US institutional sales of XRP remains.
          \end{itemize}
\end{itemize}